\setupchapter{Literaturverzeichnis}

\phantomsection % Helper, damit die Verlinkung im Inhaltsverzeichnis stimmt
\addcontentsline{toc}{section}{\sectionname}

\section*{\sectionname}
\printbibliography[heading=none] % heading=none, da Überschrift bereits durch \section* gesetzt

\begin{table}[hp]
    \centering
    \let\oldaddcontentsline\addcontentsline % Caption soll nicht im Tabellenverzeichnis erscheinen
    \renewcommand{\addcontentsline}[3]{}
    \caption{Übersicht über die verwendeten KI-basierten Werkzeuge}\label{tab:ki-werkzeuge}
    \let\addcontentsline\oldaddcontentsline
    \begin{tabularx}{\textwidth}{|l|X|}
        \hline
        Werkzeug & Beschreibung der Nutzung \\
        \hline
        ChatGPT &
        – Verständnis von Grundbegriffen im Themenfeld energiereicher Strahlungsarten (Kapitel 3.4)\newline
        – Recherche und Identifikation von Literaturstellen zu „Auswirkungen energiereicher Strahlen auf das Wachstum von Gänseblümchen“ (Kapitel 4) \\
        \hline
        ChatPDF &
        Recherche und Zusammenfassung von wissenschaftlichen Studien im Themenfeld „Messung von energiereicher Strahlung“ (Kapitel 4 und 6.2) \\
        \hline
        Microsoft CoPilot &
        – Nutzung von CoPilot in Microsoft Word 365 für Korrektur- und Formulierungshilfe (gesamt)\newline
        – Übersetzung von Textpassagen zwischen deutsch und englisch (gesamt) \\
        \hline
    \end{tabularx}
\end{table}